% Clase 25 --- El espectro electromagnetico (§6.2).
% Escala logaritmica en lambda: es la unica manera de que entren veinte ordenes
% de magnitud en un renglon. El visible se dibuja del ancho que le toca (una
% rendija) y se amplia aparte, que es justo la aclaracion que hace el docente.
\begin{center}
\begin{tikzpicture}[scale=1]

\begin{scope}[xscale=0.72]
  % ---- bandas (coordenada x = log10 de lambda en metros) ----
  \fill[figblue!45]  (-13,0) rectangle (-11,0.62);
  \fill[figblue!30]  (-11,0) rectangle (-9,0.62);
  \fill[figblue!18]  (-9,0)  rectangle (-6.4,0.62);
  \fill[figamber!35] (-6.15,0) rectangle (-3,0.62);
  \fill[figamber!22] (-3,0)  rectangle (-0.5,0.62);
  \fill[figgray!22]  (-0.5,0) rectangle (4,0.62);
  \fill[figred]      (-6.4,0) rectangle (-6.15,0.62);
  \draw[figgray, line width=0.7pt] (-13,0) rectangle (4,0.62);

  % ---- llamada al visible (va antes de los rótulos: el fondo blanco los tapa) ----
  \draw[figred, line width=0.8pt] (-6.4,0) -- (-8.6,-1.87);
  \draw[figred, line width=0.8pt] (-6.15,0) -- (2.0,-1.87);

  % ---- eje ----
  \draw[figaxis] (-13.4,-0.15) -- (4.6,-0.15);
  \foreach \xx/\lbl in {-12/\SI{1}{\pico\metre}, -9/\SI{1}{\nano\metre},
                        -3/\SI{1}{\milli\metre},
                        0/\SI{1}{\metre}, 3/\SI{1}{\kilo\metre}}{
    \draw[figgray, line width=0.6pt] (\xx,-0.1) -- (\xx,-0.28);
    \node[figlblaux, anchor=north, inner sep=2pt, fill=white] at (\xx,-0.3) {\lbl};}
  \node[figlblaux, anchor=west] at (4.7,-0.15) {$\lambda$};

  % ---- rótulos de banda ----
  \node[figlblaux, anchor=south, inner sep=2pt] at (-12,0.66) {rayos $\gamma$};
  \node[figlblaux, anchor=south, inner sep=2pt] at (-10,1.06) {rayos X};
  \node[figlblaux, anchor=south, inner sep=2pt] at (-7.7,0.66) {ultravioleta};
  \node[figlblaux, anchor=south, inner sep=2pt] at (-4.6,1.06) {infrarrojo};
  \node[figlblaux, anchor=south, inner sep=2pt] at (-1.75,0.66) {microondas};
  \node[figlblaux, anchor=south, inner sep=2pt] at (1.75,1.06)
    {radio (FM $\to$ onda larga)};

\end{scope}

  % ---- el visible ampliado (fuera del scope escalado) ----
  \fill[violet!75]  (-6.2,-1.95) rectangle (-4.55,-1.35);
  \fill[blue!65]    (-4.55,-1.95) rectangle (-2.9,-1.35);
  \fill[green!60!black] (-2.9,-1.95) rectangle (-1.25,-1.35);
  \fill[yellow!85!black] (-1.25,-1.95) rectangle (0.4,-1.35);
  \fill[orange!85]  (0.4,-1.95) rectangle (2.05,-1.35);
  \fill[red!80]     (2.05,-1.95) rectangle (3.7,-1.35);
  \draw[figgray, line width=0.7pt] (-6.2,-1.95) rectangle (3.7,-1.35);
  \node[figlblaux, anchor=north] at (-6.2,-2.0) {\SI{400}{\nano\metre}};
  \node[figlblaux, anchor=north] at (3.7,-2.0) {\SI{700}{\nano\metre}};
  \node[figlblaux, anchor=north] at (-1.25,-2.05)
    {el visible, ampliado (pico solar en el verde)};

\end{tikzpicture}\\[0.5em]
{\small\itshape La escala es logarítmica porque de otro modo no entran veinte
órdenes de magnitud en un renglón, y esa compresión es la que hace visible el
punto de fondo: \textbf{todas} estas ondas son lo mismo ---campos $\mathbf{E}$ y
$\mathbf{B}$ oscilando--- y en el vacío \textbf{todas} viajan a la misma $c$, de
modo que dar $\lambda$ o dar $\nu = c/\lambda$ es exactamente la misma
información. Lo único que las distingue es la escala espacial de la oscilación, y
de ahí salen todas las diferencias prácticas: qué las produce, qué las absorbe y
qué tamaño de antena hace falta para captarlas. Nótese lo angosto que es el
visible en la barra de arriba: es una rendija dentro de un espectro enorme, y
hubo que ampliarla aparte para poder rotularla. Las fronteras entre regiones no
son nítidas ---la misma radiación se llama rayo X o rayo $\gamma$ según cómo se
haya producido, no según su frecuencia--- y los prefijos \emph{infra} y
\emph{ultra} se entienden respecto de la \textbf{frecuencia}: el infrarrojo tiene
$\lambda$ más larga que el rojo, y por eso queda a su derecha.}
\end{center}
